\chapter{Performance Engineering}
\label{ch:performance}

\section*{Tujuan Pembelajaran}
\begin{itemize}
  \item Menetapkan anggaran performa berupa angka, yaitu p95 tiap alamat,
        ukuran jawaban, dan ambang Core Web Vitals, lalu menegakkannya sebagai
        \textit{gate} otomatis.
  \item Mengukur aplikasi di bawah beban memakai \textit{load generator},
        lalu membaca titik jenuh dari hubungan antara konkurensi,
        \textit{throughput}, dan latensi.
  \item Menemukan bagian termahal lewat \textit{profiling}, lalu membedakan
        perbaikan yang menyentuh sebabnya dari perbaikan yang hanya menyentuh
        gejalanya.
\end{itemize}

\section{Pendahuluan}

Bab 2 menetapkan \textit{baseline} berupa angka, dan Bab 3 serta Bab 4
memperbaiki dua sumber lambat yang paling sering ditemui, yaitu \textit{query} yang
boros dan perhitungan yang diulang-ulang. Yang belum ada adalah aturan yang
menyatakan kapan sebuah aplikasi dinyatakan cukup cepat. Tanpa aturan itu,
perbaikan performa berhenti ketika seseorang merasa sudah cukup, bukan ketika
angkanya tercapai.

Anggaran performa mengisi kekosongan tersebut. Anggaran adalah sekumpulan
angka yang disepakati lebih dahulu, misalnya p95 satu alamat tidak boleh
melewati 60 milidetik dan jawabannya tidak boleh melewati 8 kibibyte.
Sesudahnya, pertanyaan "apakah perubahan ini boleh digabungkan" dijawab oleh
pengukuran, bukan oleh pendapat. Bab ini memasang anggaran itu sebagai
\textit{gate} yang benar-benar memblokir, lengkap dengan status keluar yang
dapat dibaca \textit{Continuous Integration}.

Pembahasan berjalan dalam satu putaran yang berulang: menetapkan anggaran,
mengukur di bawah beban, mencari bagian termahal lewat \textit{profiling},
memperbaiki, lalu mengukur ulang. Referensi utamanya adalah ambang Core Web
Vitals dari Google \cite{webdev-vitals}, praktik anggaran performa
\cite{webdev-budget}, serta metode pengukuran sistem dari Gregg
\cite{gregg2020systems} untuk membaca hubungan antara beban, \textit{throughput},
dan waktu tunggu.

\section{Anggaran Performa}

Anggaran performa adalah daftar angka yang tidak boleh dilewati, ditulis di
satu berkas dan diperlakukan seperti kode. Isinya tiga jenis batas. Pertama,
latensi sisi server pada persentil tertentu, biasanya p95, karena rata-rata
menyembunyikan pengalaman terburuk. Kedua, ukuran jawaban di kabel, karena
byte yang dikirim menentukan waktu di jaringan yang lambat. Ketiga, ambang
Core Web Vitals, karena itulah yang benar-benar dirasakan pengguna di browser
\cite{webdev-vitals}.

Analoginya seperti anggaran belanja kepanitiaan. Angkanya ditetapkan sebelum
belanja dimulai, bukan sesudah nota menumpuk. Setiap pengeluaran baru harus
muat di dalam pagu, dan bila tidak muat, ada pengeluaran lain yang harus
dipangkas. Anggaran performa bekerja dengan cara yang sama: fitur baru boleh
menambah waktu, asalkan masih muat, atau ada bagian lain yang dipercepat
lebih dahulu.

Tabel \ref{tab:anggaran} memuat anggaran yang dipakai bab ini, apa adanya
seperti tertulis di \berkas{sources/chapter05/anggaran.json}.

\begin{table}[!htb]
  \centering
  \caption{Anggaran performa Bab 5, dinilai pada 8 \textit{worker} dengan 25
           permintaan tiap \textit{worker}}
  \label{tab:anggaran}
  \begin{tabularx}{\textwidth}{lXrr}
    \hline
    \textbf{Alamat} & \textbf{Isinya} & \textbf{p95} & \textbf{Ukuran} \\
    \hline
    \texttt{/api/katalog} & Satu halaman katalog, 20 baris & 60 ms & 8 KB \\
    \hline
    \texttt{/api/terlaris} & Agregat buku terlaris 30 hari & 400 ms & 2 KB \\
    \hline
    \texttt{/api/katalog-penuh} & Seluruh katalog, 50.000 baris & 60 ms & 8 KB \\
    \hline
    Core Web Vitals & LCP, INP, dan CLS pada persentil 75 & \multicolumn{2}{l}{2500 ms, 200 ms, 0,1} \\
    \hline
  \end{tabularx}
\end{table}

Dua hal pada anggaran sering tertukar. Pembagiannya begini. Angka lab diukur
di komputer sendiri dengan beban buatan, sehingga dapat diulang dan dipakai
sebagai \textit{gate} sebelum kode digabungkan. Angka lapangan dikumpulkan dari
kunjungan nyata, sehingga memuat perangkat dan jaringan pengguna yang
sesungguhnya. Core Web Vitals dinilai pada persentil 75 dari kunjungan nyata
\cite{webdev-vitals}, dan karena itu bab ini mengumpulkannya lewat
\textit{JavaScript} di halaman, bukan lewat alat ukur di sisi server.

Gambar \ref{fig:gate} memperlihatkan tempat anggaran bekerja di dalam alur
pengembangan.

\begin{figure}[!htb]
  \centering
  \begin{tikzpicture}[
    font=\scriptsize,
    kotak/.style={draw, rounded corners, minimum height=0.9cm, minimum width=2.2cm,
                  align=center, font=\scriptsize\bfseries},
    panah/.style={-{Stealth[length=2mm]}, thick}
  ]
    \node[kotak, fill=blue!8] (ubah) at (0,0) {Perubahan\\kode};
    \node[kotak, fill=blue!8] (ukur) at (3.1,0) {Ukur beban\\{\tiny\mdseries beban.py}};
    \node[kotak, fill=orange!12] (nilai) at (6.2,0) {Nilai anggaran\\{\tiny\mdseries anggaran.json}};
    \node[kotak, fill=green!10] (gabung) at (9.6,0.9) {Status 0\\{\tiny\mdseries boleh digabung}};
    \node[kotak, fill=red!22] (tolak) at (9.6,-0.9) {Status 1\\{\tiny\mdseries merge diblokir}};

    \draw[panah] (ubah) -- (ukur);
    \draw[panah] (ukur) -- (nilai);
    \draw[panah] (nilai.east) -- ++(0.5,0) |- (gabung.west);
    \draw[panah] (nilai.east) -- ++(0.5,0) |- (tolak.west);
    \node[font=\tiny, gray, anchor=east] at (7.7,0.72) {seluruh batas muat};
    \node[font=\tiny, gray, anchor=east] at (7.7,-0.72) {ada batas terlampaui};
  \end{tikzpicture}
  \caption{Anggaran dinilai pada tiap perubahan, dan vonisnya dibaca
           \textit{Continuous Integration} sebagai status keluar}
  \label{fig:gate}
\end{figure}

Listing \ref{lst:vitals} memperlihatkan pengukurnya. Ketiga metrik diambil
dari \texttt{PerformanceObserver} milik browser, lalu dikirim sekali saat
halaman ditinggalkan.

\begin{lstlisting}[language=TypeScript, caption={Pengukur Core Web Vitals pada kunjungan nyata, tersedia di \berkas{sources/chapter05/web-vitals.ts}}, label={lst:vitals}]
function catatLcp(daftar: PerformanceObserverEntryList): void {
  const entri = daftar.getEntries();
  const terakhir = entri[entri.length - 1];
  if (terakhir !== undefined) {
    laporan.lcp_ms = Math.round(terakhir.startTime);
  }
}

function kirimLaporan(): void {
  if (sudahDikirim) {
    return;
  }
  sudahDikirim = true;
  navigator.sendBeacon(ALAMAT_LAPORAN, JSON.stringify(laporan));
}
\end{lstlisting}

Laporan dikirim saat halaman ditinggalkan, bukan saat selesai dimuat. Alasannya
satu: INP baru diketahui setelah pengguna sempat berinteraksi, dan CLS masih
dapat bertambah sampai halaman ditutup. Pengiriman memakai
\texttt{navigator.sendBeacon}, sehingga permintaannya tetap terkirim walaupun
tabnya sudah ditutup.

\section{Metodologi Profiling}

Aturan pertama \textit{profiling} adalah mengukur lebih dahulu, baru
memperbaiki. Tanpa itu, yang diperbaiki adalah bagian yang paling mudah
ditebak, bukan bagian yang paling mahal. Aturan kedua, ukur pada tempat yang
tepat. Angka yang dirasakan klien dan angka yang terlihat di dalam server
dapat berbeda jauh, dan selisihnya justru sering menjadi petunjuk terpenting.

Listing \ref{lst:profil} memperlihatkan cara memprofil pekerjaan sisi server
tanpa melewati jaringan, memakai \texttt{cProfile} dari pustaka standar
Python \cite{pythondocs-profile}.

\begin{lstlisting}[style=PythonStyle, caption={Memprofil pekerjaan satu alamat di dalam proses, tersedia di \berkas{sources/chapter05/profil-endpoint.py}}, label={lst:profil}]
def kerjakan_katalog_penuh(pool_koneksi):
  """Mengerjakan seluruh pekerjaan alamat katalog penuh, sampai byte JSON."""
  return json.dumps(ambil_katalog_penuh(pool_koneksi)).encode()


def ulangi_pekerjaan(pekerjaan, pool_koneksi, jumlah_ulangan):
  """Menjalankan satu pekerjaan berulang kali, lalu melaporkan ukuran akhirnya.

  Diulang agar waktu tiap fungsi terkumpul cukup banyak untuk dibandingkan,
  bukan tenggelam oleh derajat ketelitian pengukur waktunya.
  """
  isi_jawaban = b""
  for _ in range(jumlah_ulangan):
    isi_jawaban = pekerjaan(pool_koneksi)
  return isi_jawaban
\end{lstlisting}

Tabel \ref{tab:profil} memuat hasilnya untuk alamat katalog penuh, 20 ulangan,
diurutkan menurut waktu kumulatif.

\begin{table}[!htb]
  \centering
  \caption{Fungsi termahal pada alamat katalog penuh, 20 ulangan, hasil
           \texttt{cProfile}}
  \label{tab:profil}
  \begin{tabularx}{\textwidth}{Xrr}
    \hline
    \textbf{Fungsi} & \textbf{Kumulatif} & \textbf{Bagian} \\
    \hline
    Seluruh pekerjaan satu alamat & 7,52 s & 100\% \\
    \hline
    \texttt{ambil\_katalog\_penuh}, \textit{query} dan pembentukan kamus & 5,10 s & 68\% \\
    \hline
    \texttt{json.dumps}, serialisasi jawaban & 1,99 s & 26\% \\
    \hline
    \texttt{baris\_menjadi\_kamus}, dipanggil 1.000.000 kali & 1,61 s & 21\% \\
    \hline
    \texttt{fetchall}, memindahkan baris ke memori & 1,28 s & 17\% \\
    \hline
  \end{tabularx}
\end{table}

Angka tersebut menjawab pertanyaan yang benar. Waktu tidak habis di \textit{query}-nya,
melainkan di memindahkan dan mengubah 50.000 baris menjadi teks. \textit{Query} yang
lebih pintar tidak akan menolong, karena yang mahal adalah jumlah barisnya.
Perbaikan yang menyentuh sebabnya hanya satu, yaitu mengirim lebih sedikit
baris.

Profil alamat agregat memberi gambaran yang berbeda. Dari 3,89 detik untuk 20
ulangan, 3,84 detik atau 99 persen dihabiskan di \texttt{connection.wait},
yaitu menunggu jawaban basis data. Di sana tidak ada yang dapat diperbaiki di
kode aplikasi. Yang tersisa adalah memperbaiki \textit{query}-nya seperti Bab 3, atau
menyimpan hasilnya seperti Bab 4.

\subsection*{Ketika klien dan server tidak sepakat}

Profil sisi server untuk alamat katalog berhalaman menunjukkan pekerjaan yang
sangat ringan, di bawah satu milidetik. Namun \textit{load generator} justru
mencatat 43,06 milidetik pada satu \textit{worker}. Selisih sebesar itu tidak
mungkin berasal dari kode yang sudah diprofil.

Penyebabnya ada di lapisan transport. Header dan isi jawaban ditulis sebagai
dua segmen TCP yang terpisah. Algoritma Nagle \cite{rfc896} menahan segmen
kecil kedua sampai segmen pertama menerima ACK (\textit{acknowledge}),
yaitu tanda terima dari penerima, sedangkan penerima menunda
ACK-nya. Keduanya bertemu pada tambahan tetap sekitar 40 milidetik untuk tiap
jawaban kecil. Tabel \ref{tab:nagle} memuat pengukuran sebelum dan sesudah
soketnya disetel dengan \texttt{TCP\_NODELAY}.

\begin{table}[!htb]
  \centering
  \caption{Pengaruh \texttt{TCP\_NODELAY}, p50 dalam milidetik, hasil
           pengukuran nyata}
  \label{tab:nagle}
  \begin{tabularx}{\textwidth}{Xrrr}
    \hline
    \textbf{Alamat dan beban} & \textbf{Sebelum} & \textbf{Sesudah} & \textbf{Rasio} \\
    \hline
    \texttt{/api/katalog}, 1 \textit{worker} & 43,06 & 0,58 & 74$\times$ \\
    \hline
    \texttt{/api/katalog}, 8 \textit{worker} & 42,99 & 2,90 & 15$\times$ \\
    \hline
    \texttt{/api/terlaris}, 1 \textit{worker} & 125,44 & 76,53 & 1,6$\times$ \\
    \hline
    \texttt{/api/katalog-penuh}, 1 \textit{worker} & 149,48 & 157,83 & 0,9$\times$ \\
    \hline
  \end{tabularx}
\end{table}

Baris terakhir sengaja dibiarkan apa adanya. Alamat katalog penuh tidak
membaik, dan memang tidak seharusnya membaik, karena jawabannya besar sehingga
segmennya tidak pernah kecil. Tambahan tetap 40 milidetik hanya terasa pada
jawaban kecil, dan justru di sanalah ia paling merusak: alamat yang
seharusnya di bawah satu milidetik menjadi puluhan kali lipat lebih lambat.

\section{Ukuran Jawaban dan Analisis Bundle}

Waktu yang dirasakan pengguna tidak hanya ditentukan kecepatan server, tetapi
juga banyaknya byte yang harus melintas. Pada jaringan 4G yang lambat, satu
megabyte tambahan berarti detik tambahan, berapa pun cepatnya basis data.
Karena itu anggaran memuat batas ukuran, bukan hanya batas waktu.

Tabel \ref{tab:ukuran-jawaban} memuat ukuran jawaban ketiga alamat, diukur
dengan \texttt{curl}, dalam dua keadaan: tanpa \textit{compress} dan dengan
\texttt{Accept-Encoding: gzip}.

\begin{table}[!htb]
  \centering
  \caption{Ukuran jawaban di kabel, hasil pengukuran nyata dengan
           \texttt{curl}}
  \label{tab:ukuran-jawaban}
  \begin{tabularx}{\textwidth}{Xrrr}
    \hline
    \textbf{Alamat} & \textbf{Tanpa gzip} & \textbf{Dengan gzip} & \textbf{Rasio} \\
    \hline
    \texttt{/api/katalog}, 20 baris & 1.545 B & 292 B & 5,3$\times$ \\
    \hline
    \texttt{/api/terlaris}, 10 baris & 557 B & 557 B & 1,0$\times$ \\
    \hline
    \texttt{/api/katalog-penuh}, 50.000 baris & 4.214.788 B & 281.905 B & 15,0$\times$ \\
    \hline
  \end{tabularx}
\end{table}

Baris tengah tampak janggal, dan memang disengaja. Jawaban 557 byte tidak
di-\textit{compress} sama sekali, karena aplikasi hanya meng-\textit{compress}
jawaban yang lebih besar dari 1.024 byte. Pada jawaban sekecil itu, gzip
menambah header dan waktu pemrosesan, sementara byte yang dihemat tidak
seberapa. Listing \ref{lst:compress} memperlihatkan aturannya.

\begin{lstlisting}[style=PythonStyle, caption={Aturan kapan jawaban di-compress, tersedia di \berkas{sources/chapter05/aplikasi.py}}, label={lst:compress}]
def compress_bila_diminta(isi_jawaban, nilai_accept_encoding):
  """Meng-compress isi jawaban bila klien menerima gzip dan isinya cukup besar.

  Mengembalikan pasangan (isi, nilai Content-Encoding atau None). Jawaban
  kecil dibiarkan apa adanya, karena gzip justru menambah byte dan waktu.
  """
  menerima_gzip = "gzip" in (nilai_accept_encoding or "")
  if not menerima_gzip or len(isi_jawaban) < BATAS_COMPRESS_BYTE:
    return isi_jawaban, None
  return gzip.compress(isi_jawaban), "gzip"
\end{lstlisting}

Baris terakhir tabel memberi pelajaran yang lebih keras. Rasio 15 kali lipat
memang mengesankan, tetapi 281.905 byte tetap 34 kali lipat lebih besar
daripada anggaran 8 kibibyte. \textit{Compress} memperkecil gejala, sedangkan
sebabnya adalah jumlah baris yang dikirim. Perbaikan yang benar memotong
jumlah barisnya, dan itulah yang dilakukan alamat katalog berhalaman.

\subsection*{Bundle di sisi browser}

Aturan yang sama berlaku untuk kode yang dikirim ke browser. Tiap kilobyte
JavaScript harus diunduh, di-\textit{parse}, lalu dijalankan, dan ketiganya
menunda interaksi pertama. Karena itu ukuran \textit{bundle} ikut dianggarkan.

Tiga cara menekan ukurannya dipakai bersama. \textit{Tree shaking} membuang
kode yang tidak pernah dipanggil saat \textit{bundle} disusun. \textit{Code
splitting} memecah \textit{bundle} menurut halaman atau fitur, sehingga
halaman katalog tidak ikut mengunduh kode halaman laporan. \textit{Lazy
loading} menunda pengunduhan potongan itu sampai benar-benar dibutuhkan.

Halaman demo bab ini memuat dua modul terpisah, dan Tabel
\ref{tab:ukuran-bundle} memuat ukurannya. Pemisahan itu sendiri sudah
merupakan bentuk paling sederhana dari \textit{code splitting}: halaman yang
tidak punya katalog cukup memuat modul pengukurnya saja.

\begin{table}[!htb]
  \centering
  \caption{Ukuran berkas halaman demo sesudah dikompilasi, hasil pengukuran
           nyata}
  \label{tab:ukuran-bundle}
  \begin{tabularx}{\textwidth}{Xrr}
    \hline
    \textbf{Berkas} & \textbf{Apa adanya} & \textbf{Dengan gzip} \\
    \hline
    \berkas{statis/index.html} & 832 B & 482 B \\
    \hline
    \berkas{statis/web-vitals.js} & 2.615 B & 1.156 B \\
    \hline
    \berkas{statis/katalog.js} & 1.591 B & 817 B \\
    \hline
    Seluruh JavaScript halaman & 4.206 B & 1.697 B \\
    \hline
  \end{tabularx}
\end{table}

Angka tersebut kecil karena halamannya memang kecil. Yang perlu dibawa bukan
angkanya, melainkan kebiasaannya: ukuran \textit{bundle} diukur pada tiap
perubahan, lalu dibandingkan dengan anggaran, persis seperti latensi.

\section{Lazy Loading serta Optimasi Gambar dan Font}

Pada halaman web kebanyakan, gambar adalah penyumbang byte terbesar, dan font
menyusul di belakangnya \cite{almanac2022page}. Keduanya juga paling mudah
diperbaiki, karena perbaikannya tidak menuntut perubahan logika aplikasi.

Empat perbaikan yang berdampak besar, ditulis berurutan dari yang paling
murah:

\begin{enumerate}
  \item Ukuran gambar disesuaikan dengan ukuran tampilnya. Gambar 2.000 piksel
        yang ditampilkan selebar 400 piksel membuang sekitar 96 persen
        pikselnya.
  \item Format modern dipakai, misalnya WebP atau AVIF, yang pada mutu setara
        menghasilkan berkas jauh lebih kecil daripada JPEG
        \cite{webdev-gambar}.
  \item Gambar di luar layar pertama ditunda dengan atribut
        \texttt{loading="lazy"}, sehingga tidak ikut berebut pita dengan
        gambar yang menentukan LCP.
  \item Font diberi \texttt{font-display: swap} agar teks tampil memakai font
        cadangan lebih dahulu, dan hanya varian yang dipakai yang diunduh.
\end{enumerate}

Satu hal penting menyertainya. Gambar yang menjadi elemen LCP justru tidak
boleh ditunda, karena penundaan itu memperlambat metrik yang sedang
diperbaiki \cite{webdev-vitals}. Untuk gambar tersebut yang dipakai adalah
kebalikannya, yaitu prioritas tinggi lewat \texttt{fetchpriority="high"}.

Ukuran gambar juga berpengaruh pada CLS. Elemen gambar tanpa atribut
\texttt{width} dan \texttt{height} membuat browser tidak dapat menyediakan
ruangnya lebih dahulu, sehingga teks di bawahnya melompat begitu gambarnya
tiba. Pergeseran itu terbaca sebagai CLS pada pengukur di Listing
\ref{lst:vitals}, dan cara mencegahnya cukup mencantumkan ukurannya.

\section{Streaming Server-Side Rendering}

Pada halaman yang dirakit di server, seluruh halaman biasanya dikirim setelah
seluruh datanya siap. Akibatnya pengguna menatap layar kosong selama server
bekerja. \textit{Streaming} mengubah urutannya: bagian yang sudah siap
dikirim lebih dahulu, sedangkan bagian yang masih dihitung menyusul dalam
potongan berikutnya. Mekanismenya adalah \textit{chunked transfer encoding}
milik HTTP/1.1 \cite{rfc9112}.

Analoginya seperti menyajikan makanan di kantin. Cara pertama menunggu seluruh
pesanan matang, baru semuanya diantar sekaligus. Cara kedua mengantar yang
sudah matang lebih dahulu. Total waktu memasak sama, tetapi yang menunggu
merasakan keduanya sangat berbeda.

Listing \ref{lst:streaming} memperlihatkan penerapannya. Header dan potongan
pertama dikirim sebelum \textit{query} agregat dijalankan.

\begin{lstlisting}[style=PythonStyle, caption={Mengirim header lebih dahulu, baru menghitung, tersedia di \berkas{sources/chapter05/aplikasi.py}}, label={lst:streaming}]
  def layani_terlaris_streaming(self):
    """Mengirim header lebih dahulu, baru menghitung agregatnya.

    Bagian yang sudah siap dikirim duluan, sehingga waktu sampai byte pertama
    tidak menunggu query selesai. Waktu total tidak berubah.
    """
    self.send_response(200)
    self.send_header("Content-Type", "application/json")
    self.send_header("Cache-Control", CACHE_CONTROL_API)
    self.send_header("Transfer-Encoding", "chunked")
    self.end_headers()
    self.kirim_potongan(b'{"item":[')
    daftar_terlaris = jalankan_query_terlaris(self.server.pool_koneksi)
\end{lstlisting}

Tabel \ref{tab:streaming} memuat hasil pengukurannya dengan \texttt{curl},
tiga permintaan berurutan untuk tiap alamat.

\begin{table}[!htb]
  \centering
  \caption{Waktu sampai byte pertama dan waktu total, dalam milidetik, hasil
           pengukuran nyata}
  \label{tab:streaming}
  \begin{tabularx}{\textwidth}{Xrr}
    \hline
    \textbf{Alamat} & \textbf{Byte pertama} & \textbf{Total} \\
    \hline
    \texttt{/api/terlaris} & 118 sampai 196 & 119 sampai 196 \\
    \hline
    \texttt{/api/terlaris-streaming} & 0,7 sampai 0,9 & 95 sampai 103 \\
    \hline
  \end{tabularx}
\end{table}

Dua hal terbaca dari tabel tersebut. Pertama, waktu sampai byte pertama turun
dari ratusan milidetik menjadi kurang dari satu milidetik, karena browser
menerima header dan potongan pertama sebelum \textit{query} dimulai. Kedua, waktu
totalnya tidak membaik, dan memang tidak seharusnya membaik. \textit{Streaming}
memindahkan waktu tunggu, bukan menghapusnya. Yang diperbaiki adalah
pengalaman menunggu, sedangkan pekerjaannya tetap sama.

\section{Tuning Basis Data dan Query}

Bab 3 menyetel \textit{query} lewat indeks, dan Bab 4 menghindari perhitungan berulang
lewat \textit{cache}. Bab ini menambahkan satu penyetelan yang sering
terlewat, yaitu memutuskan berapa banyak baris yang pantas dikirim dalam satu
jawaban.

Tabel \ref{tab:katalog} membandingkan dua alamat yang menjawab kebutuhan yang
sama, yaitu menampilkan katalog. Yang pertama mengirim satu halaman berisi 20
baris, yang kedua mengirim seluruh 50.000 baris.

\begin{table}[!htb]
  \centering
  \caption{Katalog berhalaman lawan katalog penuh, waktu dalam milidetik,
           hasil pengukuran nyata pada 8 \textit{worker}}
  \label{tab:katalog}
  \begin{tabularx}{\textwidth}{Xrrr}
    \hline
    \textbf{Ukuran} & \textbf{Berhalaman} & \textbf{Penuh} & \textbf{Rasio} \\
    \hline
    p50 & 9,24 & 2.202,68 & 238$\times$ \\
    \hline
    p95 & 14,02 & 2.998,55 & 214$\times$ \\
    \hline
    \textit{Throughput}, permintaan per detik & 811,8 & 3,5 & 232$\times$ \\
    \hline
    Ukuran jawaban di kabel & 0,3 KB & 275,3 KB & 918$\times$ \\
    \hline
  \end{tabularx}
\end{table}

Selisih ratusan kali lipat itu tidak berasal dari \textit{query} yang lebih pintar.
Kedua alamat memakai tabel dan indeks yang sama. Yang berbeda hanya jumlah
baris yang diambil, diubah menjadi kamus, lalu di-\textit{serialize}. Profil
pada Tabel \ref{tab:profil} sudah menunjukkannya: \texttt{json.dumps} dan
\texttt{baris\_menjadi\_kamus} menghabiskan hampir separuh waktu, dan keduanya
tumbuh lurus mengikuti jumlah baris.

Aturan yang dapat dibawa keluar dari sini ada dua. Pertama, setiap alamat
yang mengembalikan daftar wajib berhalaman, dan ukuran halamannya dibatasi di
sisi server, bukan dipercayakan kepada klien. Listing
\ref{lst:halaman} memperlihatkan pembatasan itu. Kedua, kolom yang tidak
dipakai tidak ikut diambil, karena tiap kolom menambah byte yang diubah dan
dikirim.

\begin{lstlisting}[style=PythonStyle, caption={Batas ukuran halaman ditegakkan di sisi server, tersedia di \berkas{sources/chapter05/aplikasi.py}}, label={lst:halaman}]
def baca_parameter_halaman(query_alamat):
  """Membaca parameter halaman dan ukuran dari alamat, beserta batasnya.

  Ukuran dibatasi agar satu permintaan tidak dapat meminta seluruh tabel
  hanya dengan mengubah alamat.
  """
  parameter = parse_qs(query_alamat)
  halaman = max(1, int(parameter.get("halaman", ["1"])[0]))
  ukuran = int(parameter.get("ukuran", [str(UKURAN_HALAMAN_BAWAAN)])[0])
  return halaman, min(max(1, ukuran), UKURAN_HALAMAN_MAKSIMUM)
\end{lstlisting}

Alamat agregat menempuh jalan yang berbeda. Jumlah barisnya sudah kecil,
hanya sepuluh, tetapi p95-nya 930,99 milidetik pada 8 \textit{worker},
sehingga tetap melewati anggaran 400 milidetik. Profilnya menunjukkan 99
persen waktu dihabiskan menunggu basis data, jadi memperbaiki kode aplikasi
tidak akan menolong. Yang tersisa adalah menyimpan hasilnya seperti Bab 4,
atau menambah kapasitas basis datanya.

\section{Load Testing dan Identifikasi Bottleneck}

Mengukur satu permintaan pada saat sepi menjawab pertanyaan yang salah.
Pengguna datang bersamaan, dan perilaku sistem saat ramai tidak dapat
ditebak dari perilakunya saat sepi. \textit{Load testing} memberi beban yang
terkendali, lalu mencatat apa yang terjadi pada latensi dan
\textit{throughput}.

Listing \ref{lst:beban} memperlihatkan inti \textit{load generator} yang
dipakai bab ini. Bentuknya disebut \textit{closed loop}: tiap \textit{worker}
menunggu jawaban sebelum mengirim permintaan berikutnya, persis seperti satu
pengguna yang sabar.

\begin{lstlisting}[style=PythonStyle, caption={Satu worker mengirim permintaan berurutan lewat koneksi yang dipakai ulang, tersedia di \berkas{sources/chapter05/beban.py}}, label={lst:beban}]
def jalankan_satu_worker(jalur, jumlah_permintaan, catatan):
  """Mengirim permintaan berurutan lewat satu koneksi, lalu mencatat hasilnya.

  Koneksi dibuka sekali dan dipakai ulang, sehingga yang terukur adalah waktu
  server, bukan waktu membuka koneksi baru.
  """
  koneksi = HTTPConnection(HOST_BAWAAN, PORT_BAWAAN, timeout=BATAS_WAKTU_DETIK)
  header = {"Accept-Encoding": "gzip"}
  for _ in range(jumlah_permintaan):
    waktu_mulai = time.perf_counter()
\end{lstlisting}

Tabel \ref{tab:ramp} memuat hasil menaikkan jumlah \textit{worker} pada alamat
agregat, dari satu sampai 32, dengan 15 permintaan tiap \textit{worker}.

\begin{table}[!htb]
  \centering
  \caption{Pengaruh konkurensi pada latensi dan \textit{throughput} alamat
           agregat, hasil pengukuran nyata}
  \label{tab:ramp}
  \begin{tabularx}{\textwidth}{rrrrX}
    \hline
    \textbf{Worker} & \textbf{p50} & \textbf{p95} & \textbf{Per detik} & \textbf{Yang terjadi} \\
    \hline
    1 & 156,83 & 214,61 & 6,0 & Tidak ada antrean \\
    \hline
    2 & 172,61 & 184,05 & 11,6 & Kapasitas masih terpakai \\
    \hline
    4 & 302,01 & 409,19 & 12,2 & Mulai jenuh \\
    \hline
    8 & 514,98 & 835,65 & 13,6 & Tambahan beban menjadi antrean \\
    \hline
    16 & 1.210,84 & 2.153,35 & 11,6 & \textit{Throughput} turun \\
    \hline
    32 & 2.794,87 & 3.945,97 & 11,0 & Latensi meledak, hasil tetap \\
    \hline
  \end{tabularx}
\end{table}

Pola pada tabel tersebut adalah pola yang paling sering ditemui pada sistem
yang jenuh. Sampai dua \textit{worker}, menambah beban menambah hasil.
Sesudahnya \textit{throughput} mentok di kisaran 12 sampai 14 permintaan per
detik, sedangkan latensi terus naik hampir lurus mengikuti jumlah
\textit{worker}. Pada 32 \textit{worker}, p95 menjadi 18 kali lipat p95 pada
satu \textit{worker}, sementara hasilnya justru sedikit lebih rendah. Gambar \ref{fig:jenuh}
menggambarkan keduanya dalam satu sumbu.

\begin{figure}[!htb]
  \centering
  \begin{tikzpicture}[font=\scriptsize]
    \begin{axis}[
      width=0.86\textwidth, height=5.4cm,
      xmode=log, log basis x=2, xmin=0.85, xmax=38,
      xtick={1,2,4,8,16,32}, xticklabels={1,2,4,8,16,32},
      xlabel={Jumlah \textit{worker}}, ylabel={Permintaan per detik},
      ymin=0, ymax=16, axis y line*=left, axis x line*=bottom,
    ]
      \addplot[praditagreen, very thick, mark=*, mark size=1.6]
        coordinates {(1,6.0) (2,11.6) (4,12.2) (8,13.6) (16,11.6) (32,11.0)};
    \end{axis}
    \begin{axis}[
      width=0.86\textwidth, height=5.4cm,
      xmode=log, log basis x=2, xmin=0.85, xmax=38,
      axis y line*=right, axis x line=none, xtick=\empty,
      ylabel={p95 (milidetik)}, ymin=0, ymax=4200,
    ]
      \addplot[red!65!black, very thick, dashed, mark=square*, mark size=1.4]
        coordinates {(1,214.61) (2,184.05) (4,409.19) (8,835.65)
                     (16,2153.35) (32,3945.97)};
    \end{axis}
    \node[font=\tiny, praditagreen, anchor=west] at (1.1,3.3) {permintaan per detik};
    \node[font=\tiny, red!65!black, anchor=west] at (5.0,1.2) {p95};
  \end{tikzpicture}
  \caption{Sesudah titik jenuh, menambah \textit{worker} hanya menaikkan p95
           tanpa menambah permintaan per detik}
  \label{fig:jenuh}
\end{figure}

Titik itu disebut titik jenuh, dan membacanya memberi tiga kesimpulan
sekaligus. Kapasitas sistem ini sekitar 13 permintaan per detik untuk alamat
tersebut. Menambah \textit{worker} di atas titik jenuh tidak menambah
kapasitas, hanya memindahkan waktu tunggu ke dalam antrean. Perbaikan yang
benar harus menaikkan kapasitasnya, misalnya lewat \textit{cache} Bab 4,
bukan menaikkan konkurensi.

Hubungan ketiga angka itu dapat diperiksa dengan hukum Little
\cite{gregg2020systems}. Jumlah permintaan yang sedang dikerjakan sama dengan
\textit{throughput} dikalikan waktu tinggalnya. Pada baris 8 \textit{worker},
13,6 permintaan per detik dikalikan 0,515 detik menghasilkan 7,0, dan angka
itu memang mendekati 8. Kecocokan tersebut menjadi pemeriksaan kewarasan:
bila ketiganya tidak konsisten, yang salah biasanya pengukurannya.

\subsection*{Menjaga pengukuran tetap jujur}

Pengukuran beban mudah tercemar oleh pengukuran sebelumnya. Alamat katalog
penuh membaca seluruh tabel ratusan kali, sehingga halaman tabel lain
terdorong keluar dari \textit{shared buffers} PostgreSQL. Pengukuran alamat
agregat yang dijalankan tepat sesudahnya mencatat p95 sekitar 1.000
milidetik, padahal pengukuran yang sama dengan pemanasan lebih dahulu
mencatat 831 sampai 931 milidetik.

Karena itu \textit{gate} anggaran memakai tiga aturan. Tiap alamat dipanaskan
dengan 30 permintaan sebelum diukur. Ada jeda tiga detik sebelum tiap
pengukuran agar sisa beban sebelumnya reda. Alamat terberat diukur paling
akhir, sehingga tidak mencemari alamat lain.

\subsection*{Gate yang benar-benar memblokir}

Anggaran baru berarti bila ada yang menegakkannya. Listing \ref{lst:gate}
memperlihatkan bagian penutup \textit{gate}-nya: begitu ada satu anggaran yang
terlampaui, prosesnya keluar dengan status 1, dan \textit{Continuous
Integration} membaca status itu sebagai kegagalan.

\begin{lstlisting}[style=PythonStyle, caption={Status keluar yang memblokir merge, tersedia di \berkas{sources/chapter05/gate-anggaran.py}}, label={lst:gate}]
  jumlah_lewat = sum(1 for baris in daftar_baris if not baris[-1])
  if jumlah_lewat:
    print(f"\nGagal: {jumlah_lewat} anggaran terlampaui.")
    sys.exit(1)
  print("\nLolos: seluruh anggaran terpenuhi.")
\end{lstlisting}

Listing \ref{lst:gate-jalan} memuat keluarannya apa adanya pada aplikasi
bab ini.

\begin{lstlisting}[language=bash, caption={Menjalankan \textit{gate} anggaran, tersedia di \berkas{sources/chapter05/gate-anggaran.py}}, label={lst:gate-jalan}]
source ../.venv/bin/activate
python gate-anggaran.py

# Keluarannya:
#   Ukuran                              Anggaran           Terukur  Vonis
#   /api/katalog?halaman=1&ukuran=20 p95    60 ms         12.59 ms  lolos
#   /api/katalog?halaman=1&ukuran=20 ukuran  8 KB          0.29 KB  lolos
#   /api/terlaris p95                      400 ms        830.95 ms  LEWAT
#   /api/terlaris ukuran                     2 KB          0.54 KB  lolos
#   /api/katalog-penuh p95                  60 ms       2262.70 ms  LEWAT
#   /api/katalog-penuh ukuran                8 KB        275.30 KB  LEWAT
#   lcp_ms p75                            2500 ms   tidak ada data  LEWAT
#   inp_ms p75                             200 ms   tidak ada data  LEWAT
#   cls p75                                0.1      tidak ada data  LEWAT
#
# Gagal: 6 anggaran terlampaui.
\end{lstlisting}

Tiga baris terakhir layak diperhatikan. Core Web Vitals tidak dinyatakan
lolos hanya karena belum ada laporannya. Selama tidak ada kunjungan nyata
yang terekam, \textit{gate} menilainya gagal. Aturan itu disengaja, karena
perbaikan tanpa bukti pengukuran tidak dapat dihitung sebagai perbaikan.

\section{Ringkasan}

Performa berhenti menjadi perdebatan ketika angkanya ditulis lebih dahulu.
Anggaran performa berisi tiga jenis batas, yaitu latensi pada persentil
tertentu, ukuran jawaban di kabel, dan ambang Core Web Vitals. Ketiganya
disimpan di satu berkas, diperlakukan seperti kode, dan dinilai oleh
\textit{gate} yang keluar dengan status 1 begitu ada satu batas yang
terlampaui.

Perbaikan dimulai dari pengukuran, bukan dari tebakan. Profil alamat katalog
penuh menunjukkan waktu habis di serialisasi dan pembentukan kamus untuk
50.000 baris, sehingga \textit{query} yang lebih pintar tidak akan menolong. Profil
alamat agregat menunjukkan 99 persen waktu dihabiskan menunggu basis data,
sehingga yang perlu diperbaiki ada di luar kode aplikasi. Dua profil, dua
kesimpulan yang sama sekali berbeda.

Tempat mengukur ikut menentukan kesimpulan. Pekerjaan sisi server alamat
katalog berada di bawah satu milidetik, tetapi klien mencatat 43,06
milidetik. Selisihnya bukan berasal dari kode, melainkan dari algoritma Nagle
yang menahan segmen kecil kedua. Setelah soketnya disetel, p50 turun menjadi
0,58 milidetik, yaitu 74 kali lipat lebih cepat, tanpa satu baris pun logika
aplikasi berubah.

Beban mengubah segalanya. Alamat agregat melayani 6,0 permintaan per detik
pada satu \textit{worker} dan 13,6 pada delapan, lalu berhenti bertambah.
Menaikkan konkurensi menjadi 32 hanya menaikkan p95 dari 214,61 menjadi
3.945,97 milidetik, sedangkan hasilnya tidak bertambah. Titik jenuh itulah
kapasitas sistem, dan angka itu hanya terlihat lewat \textit{load testing},
bukan lewat pengukuran satu permintaan.

Sebagian perbaikan memindahkan waktu tunggu, bukan menghapusnya, dan itu
tetap berharga. \textit{Streaming} menurunkan waktu sampai byte pertama dari
ratusan milidetik menjadi kurang dari satu milidetik tanpa mengubah waktu
total. Yang membedakan perbaikan sungguhan dari perbaikan semu bukan
besarnya angka, melainkan adanya bukti pengukuran sebelum dan sesudahnya.

\section{Latihan}

Seluruh berkas yang dipakai pada latihan ini tersedia di
\berkas{sources/chapter05/}, dihitung relatif terhadap akar repositori.
Latihan membutuhkan Docker dengan Compose versi 2, Python 3, dan Node.js
untuk mengompilasi berkas TypeScript. Listing penyiapan masuk ke direktori
itu lebih dahulu, dan seluruh perintah sesudahnya dijalankan dari sana. Bila
membuka terminal baru, masuk kembali dengan \texttt{cd sources/chapter05}
dari akar repositori.

Peringatan yang sama seperti bab sebelumnya berlaku di sini.
\berkas{sources/chapter05/beban.py} dan
\berkas{sources/chapter05/gate-anggaran.py} sengaja membebani, sehingga
keduanya hanya diarahkan ke aplikasi milik sendiri. Hasil tiap mahasiswa akan
berbeda, karena angkanya bergantung pada perangkat keras, versi PostgreSQL,
dan beban lain yang sedang berjalan di komputer yang sama. Yang dilaporkan
adalah pola perbandingannya beserta penjelasannya, bukan angka persisnya.

Penyiapannya dilakukan sekali, seperti pada Listing \ref{lst:c5-siapkan}.

\begin{lstlisting}[language=bash, caption={Menyiapkan basis data, data uji, dan berkas statis, memakai \berkas{sources/chapter05/siapkan-data.sh}}, label={lst:c5-siapkan}]
cd sources/chapter05
docker compose up -d
./siapkan-data.sh
python3 -m venv ../.venv
../.venv/bin/pip install "psycopg[binary,pool]" sqlalchemy
npx -p typescript tsc --target es2022 --module es2022 --strict \
  --lib es2022,dom --outDir statis web-vitals.ts katalog.ts
\end{lstlisting}

Aplikasinya dinyalakan di satu terminal dan dibiarkan menyala selama latihan
berlangsung, seperti pada Listing \ref{lst:c5-aplikasi}.

\begin{lstlisting}[language=bash, caption={Menyalakan aplikasi yang diukur, tersedia di \berkas{sources/chapter05/aplikasi.py}}, label={lst:c5-aplikasi}]
source ../.venv/bin/activate
python aplikasi.py

# Keluarannya:
# Melayani di http://127.0.0.1:8010/  (Ctrl+C untuk berhenti)
\end{lstlisting}

\subsection*{Latihan 1: Menegakkan Anggaran Performa}

Latihan ini membuktikan bahwa anggaran hanya berguna bila ada yang
menegakkannya, dan bahwa vonisnya dapat dibaca mesin.

\begin{enumerate}
  \item Baca \berkas{sources/chapter05/anggaran.json}, lalu tulis dugaan
        untuk tiap alamat: lolos atau tidak, dan mengapa. Dugaan ditulis
        lebih dahulu agar tidak menyesuaikan diri dengan hasilnya.
  \item Jalankan \textit{gate}-nya seperti pada Listing \ref{lst:c5-lat1-jalan},
        lalu catat vonis tiap baris ke Tabel \ref{tab:c5-lat1-hasil}. Tabel
        \ref{tab:c5-lat1-contoh} memperlihatkan contoh yang sudah terisi.
  \item Catat status keluar prosesnya dengan \texttt{echo \$?} tepat sesudah
        \textit{gate} selesai. Angka itulah yang dibaca \textit{Continuous
        Integration}.
  \item Ubah satu angka anggaran menjadi longgar, misalnya p95 alamat agregat
        menjadi 2000 milidetik, lalu jalankan ulang \textit{gate}-nya dan catat
        status keluarnya.
  \item Jawab tiga pertanyaan berikut. Pertama, alamat mana yang gagal, dan
        apakah kegagalannya berasal dari waktu atau dari ukuran? Kedua,
        mengapa baris Core Web Vitals dinyatakan gagal padahal aplikasinya
        berjalan normal? Ketiga, apa akibatnya bila \textit{gate} dibuat selalu
        keluar dengan status 0?
\end{enumerate}

\begin{lstlisting}[language=bash, caption={Menjalankan \textit{gate} anggaran dan membaca status keluarnya, tersedia di \berkas{sources/chapter05/gate-anggaran.py}}, label={lst:c5-lat1-jalan}]
source ../.venv/bin/activate
python gate-anggaran.py
echo $?

# Keluarannya:
# Gagal: 6 anggaran terlampaui.
# 1
\end{lstlisting}

\begin{table}[!htb]
  \centering
  \caption{Contoh hasil Latihan 1 yang sudah terisi, waktu dalam milidetik}
  \label{tab:c5-lat1-contoh}
  \begin{tabularx}{\textwidth}{Xrrl}
    \hline
    \textbf{Ukuran} & \textbf{Anggaran} & \textbf{Terukur} & \textbf{Vonis} \\
    \hline
    \texttt{/api/katalog} p95 & 60 & 12,59 & lolos \\
    \hline
    \texttt{/api/katalog} ukuran & 8 KB & 0,29 KB & lolos \\
    \hline
    \texttt{/api/terlaris} p95 & 400 & 830,95 & LEWAT \\
    \hline
    \texttt{/api/katalog-penuh} p95 & 60 & 2.262,70 & LEWAT \\
    \hline
    \texttt{/api/katalog-penuh} ukuran & 8 KB & 275,30 KB & LEWAT \\
    \hline
    LCP p75 & 2500 & tidak ada data & LEWAT \\
    \hline
  \end{tabularx}
\end{table}

\begin{table}[!htb]
  \centering
  \caption{Lembar isian hasil pengukuran Latihan 1}
  \label{tab:c5-lat1-hasil}
  \begin{tabularx}{\textwidth}{Xrrl}
    \hline
    \textbf{Ukuran} & \textbf{Anggaran} & \textbf{Terukur} & \textbf{Vonis} \\
    \hline
    \texttt{/api/katalog} p95 & \dots & \dots & \dots \\
    \hline
    \texttt{/api/katalog} ukuran & \dots & \dots & \dots \\
    \hline
    \texttt{/api/terlaris} p95 & \dots & \dots & \dots \\
    \hline
    \texttt{/api/katalog-penuh} p95 & \dots & \dots & \dots \\
    \hline
    \texttt{/api/katalog-penuh} ukuran & \dots & \dots & \dots \\
    \hline
    Status keluar sebelum dan sesudah anggaran dilonggarkan & \dots & \dots & \dots \\
    \hline
  \end{tabularx}
\end{table}

\subsection*{Latihan 2: Menemukan Bagian Termahal}

Latihan ini membuktikan bahwa dua alamat yang sama-sama lambat dapat
membutuhkan perbaikan yang sama sekali berbeda.

\begin{enumerate}
  \item Jalankan profil kedua alamat seperti pada Listing
        \ref{lst:c5-lat2-jalan}.
  \item Untuk tiap alamat, catat tiga fungsi teratas menurut waktu kumulatif
        beserta bagiannya terhadap waktu total, lalu isi Tabel
        \ref{tab:c5-lat2-hasil}. Tabel \ref{tab:c5-lat2-contoh}
        memperlihatkan contoh yang sudah terisi.
  \item Untuk tiap alamat, tulis satu kalimat perbaikan yang menyentuh
        sebabnya, bukan gejalanya.
  \item Jawab tiga pertanyaan berikut. Pertama, mengapa mempercepat \textit{query}
        tidak banyak menolong alamat katalog penuh? Kedua, mengapa
        mempercepat kode Python tidak menolong alamat agregat? Ketiga,
        berapa kali \texttt{baris\_menjadi\_kamus} dipanggil untuk 20
        ulangan, dan dari mana angka itu berasal?
\end{enumerate}

\begin{lstlisting}[language=bash, caption={Memprofil dua alamat, tersedia di \berkas{sources/chapter05/profil-endpoint.py}}, label={lst:c5-lat2-jalan}]
source ../.venv/bin/activate
python profil-endpoint.py katalog-penuh 20
python profil-endpoint.py terlaris 20
\end{lstlisting}

\begin{table}[!htb]
  \centering
  \caption{Contoh hasil Latihan 2 yang sudah terisi, 20 ulangan}
  \label{tab:c5-lat2-contoh}
  \begin{tabularx}{\textwidth}{lXrr}
    \hline
    \textbf{Alamat} & \textbf{Fungsi teratas} & \textbf{Kumulatif} & \textbf{Bagian} \\
    \hline
    katalog penuh & \texttt{ambil\_katalog\_penuh} & 5,10 s & 68\% \\
    \hline
    katalog penuh & \texttt{json.dumps} & 1,99 s & 26\% \\
    \hline
    katalog penuh & \texttt{baris\_menjadi\_kamus} & 1,61 s & 21\% \\
    \hline
    agregat & \texttt{connection.wait} & 3,84 s & 99\% \\
    \hline
  \end{tabularx}
\end{table}

\begin{table}[!htb]
  \centering
  \caption{Lembar isian hasil pengukuran Latihan 2}
  \label{tab:c5-lat2-hasil}
  \begin{tabularx}{\textwidth}{lXrr}
    \hline
    \textbf{Alamat} & \textbf{Fungsi teratas} & \textbf{Kumulatif} & \textbf{Bagian} \\
    \hline
    katalog penuh & \dots & \dots & \dots \\
    \hline
    katalog penuh & \dots & \dots & \dots \\
    \hline
    katalog penuh & \dots & \dots & \dots \\
    \hline
    agregat & \dots & \dots & \dots \\
    \hline
  \end{tabularx}
\end{table}

\subsection*{Latihan 3: Menemukan Titik Jenuh}

Latihan ini mencari kapasitas satu alamat, yaitu jumlah permintaan per detik
yang masih dapat dilayani sebelum latensinya meledak.

\begin{enumerate}
  \item Panaskan alamatnya dengan 30 permintaan pada satu \textit{worker},
        agar \textit{cache} basis data hangat sebelum pengukuran dimulai.
  \item Jalankan beban untuk 1, 2, 4, 8, 16, dan 32 \textit{worker} seperti
        pada Listing \ref{lst:c5-lat3-jalan}, lalu isi Tabel
        \ref{tab:c5-lat3-hasil}. Tabel \ref{tab:c5-lat3-contoh}
        memperlihatkan contoh yang sudah terisi.
  \item Tandai baris tempat \textit{throughput} berhenti bertambah. Itulah
        titik jenuhnya.
  \item Periksa hukum Little pada satu baris pilihan: kalikan
        \textit{throughput} dengan p50 dalam detik, lalu bandingkan hasilnya
        dengan jumlah \textit{worker}.
  \item Jawab tiga pertanyaan berikut. Pertama, pada jumlah \textit{worker}
        berapa \textit{throughput} berhenti bertambah, dan apa yang terjadi
        pada p95 sesudah titik itu? Kedua, mengapa menambah \textit{worker}
        di atas titik jenuh tidak menambah kapasitas? Ketiga, perbaikan mana
        yang menaikkan kapasitas, dan mana yang hanya memindahkan antrean?
\end{enumerate}

\begin{lstlisting}[language=bash, caption={Menaikkan konkurensi bertahap, tersedia di \berkas{sources/chapter05/beban.py}}, label={lst:c5-lat3-jalan}]
source ../.venv/bin/activate
python beban.py /api/terlaris 1 30
python beban.py /api/terlaris 1 15
python beban.py /api/terlaris 8 15
python beban.py /api/terlaris 32 15
\end{lstlisting}

\begin{table}[!htb]
  \centering
  \caption{Contoh hasil Latihan 3 yang sudah terisi, waktu dalam milidetik}
  \label{tab:c5-lat3-contoh}
  \begin{tabularx}{\textwidth}{rrrrX}
    \hline
    \textbf{Worker} & \textbf{p50} & \textbf{p95} & \textbf{Per detik} & \textbf{Catatan} \\
    \hline
    1 & 156,83 & 214,61 & 6,0 & Tanpa antrean \\
    \hline
    8 & 514,98 & 835,65 & 13,6 & Sudah jenuh \\
    \hline
    32 & 2.794,87 & 3.945,97 & 11,0 & Latensi meledak \\
    \hline
  \end{tabularx}
\end{table}

\begin{table}[!htb]
  \centering
  \caption{Lembar isian hasil pengukuran Latihan 3}
  \label{tab:c5-lat3-hasil}
  \begin{tabularx}{\textwidth}{rrrrX}
    \hline
    \textbf{Worker} & \textbf{p50} & \textbf{p95} & \textbf{Per detik} & \textbf{Catatan} \\
    \hline
    1 & \dots & \dots & \dots & \dots \\
    \hline
    2 & \dots & \dots & \dots & \dots \\
    \hline
    4 & \dots & \dots & \dots & \dots \\
    \hline
    8 & \dots & \dots & \dots & \dots \\
    \hline
    16 & \dots & \dots & \dots & \dots \\
    \hline
    32 & \dots & \dots & \dots & \dots \\
    \hline
  \end{tabularx}
\end{table}

\subsection*{Latihan 4: Mengumpulkan Core Web Vitals Lapangan}

Latihan ini melengkapi angka lab dengan angka lapangan, yaitu Core Web Vitals
dari kunjungan yang sungguhan.

\begin{enumerate}
  \item Pastikan berkas TypeScript sudah dikompilasi pada langkah penyiapan,
        lalu buka \texttt{http://localhost:8010/} di browser.
  \item Tekan tombol Muat katalog, tunggu daftarnya tampil, lalu tutup tab
        atau pindah ke tab lain. Laporan dikirim saat halaman ditinggalkan,
        bukan saat selesai dimuat.
  \item Periksa isi \berkas{sources/chapter05/vitals.jsonl} seperti pada
        Listing \ref{lst:c5-lat4-jalan}, lalu catat ketiga metriknya ke Tabel
        \ref{tab:c5-lat4-hasil}. Tabel \ref{tab:c5-lat4-contoh}
        memperlihatkan bentuk isian yang diharapkan.
  \item Ulangi kunjungan sebanyak lima kali, lalu jalankan ulang \textit{gate}
        anggaran dan catat perubahan vonis ketiga baris Core Web Vitals.
  \item Jawab tiga pertanyaan berikut. Pertama, mengapa laporan dikirim saat
        halaman ditinggalkan? Kedua, mengapa Core Web Vitals dinilai pada
        persentil 75 dari banyak kunjungan, bukan dari satu kunjungan?
        Ketiga, apa yang terjadi pada CLS bila gambar dipasang tanpa atribut
        \texttt{width} dan \texttt{height}?
\end{enumerate}

\begin{lstlisting}[language=bash, caption={Memeriksa laporan Core Web Vitals yang masuk, dihasilkan \berkas{sources/chapter05/web-vitals.ts}}, label={lst:c5-lat4-jalan}]
cat vitals.jsonl
source ../.venv/bin/activate
python gate-anggaran.py
\end{lstlisting}

\begin{table}[!htb]
  \centering
  \caption{Bentuk isian Latihan 4 yang diharapkan}
  \label{tab:c5-lat4-contoh}
  \begin{tabularx}{\textwidth}{lrrX}
    \hline
    \textbf{Kunjungan} & \textbf{LCP} & \textbf{INP} & \textbf{CLS} \\
    \hline
    1 & 1.840 ms & 24 ms & 0,0021 \\
    \hline
    p75 dari lima kunjungan & \dots & \dots & \dots \\
    \hline
  \end{tabularx}
\end{table}

\begin{table}[!htb]
  \centering
  \caption{Lembar isian hasil pengukuran Latihan 4}
  \label{tab:c5-lat4-hasil}
  \begin{tabularx}{\textwidth}{lrrX}
    \hline
    \textbf{Kunjungan} & \textbf{LCP} & \textbf{INP} & \textbf{CLS} \\
    \hline
    1 & \dots & \dots & \dots \\
    \hline
    2 & \dots & \dots & \dots \\
    \hline
    3 & \dots & \dots & \dots \\
    \hline
    4 & \dots & \dots & \dots \\
    \hline
    5 & \dots & \dots & \dots \\
    \hline
    p75 dari lima kunjungan & \dots & \dots & \dots \\
    \hline
  \end{tabularx}
\end{table}

Setelah seluruh latihan selesai, container dihentikan dengan
\texttt{docker compose down -v}. Opsi \texttt{-v} ikut menghapus volume
datanya, sehingga penyiapan berikutnya dimulai dari keadaan bersih.
